% problem-000243 4 分光光度法测平衡常数
\section*{4. 分光光度法测平衡常数}
分光光度法灵敏度高，可用来测定有色弱酸弱碱的解离常数或配合物的稳定常数。(本题需要修约有效数字)

(a)
（图：）

（图：）
图 4.1 (a) HIn-的结构式；(b) 紫外-可见吸收光谱

\subsection*{4-1}
溴百里酚蓝一钠盐（用NaHIn表示）是一种常用的酸碱指示剂，其结构如图4.1a所示，pH变色范围为6.0（黄色）\~7.8（蓝色）。其在水中解离度很小，用一般的分析方法很难准确测定其解离常数 $K _ { \mathfrak { d } }$ 值，但利用分光光度法可以准确、方便地测量。NaHIn的黄色和蓝色溶液的紫外-可见吸收光谱如图4.1b所示。

\subsection*{4-1}
-1 画出 HIn-共轭碱的结构式。

\subsection*{4-1}
-2NaHIn的中性溶液呈绿色，简要解释原因。

\subsection*{4-1}
\~3图4.1b中实线对应的是什么颜色溶液的吸收光谱，简要解释原因。

\subsection*{4-1}
-4 取3 份等量的 NaHIn 溶液，分别准确添加相同体积的 HCl、NaOH与 $\mathsf { p H } = 7 . 0 0$ 的缓冲溶液，使溶液呈黄色、蓝色、绿色。室温下选择不同波长测量各溶液的吸光度，结果见表4.1。推导NaHIn的 $K _ { \mathfrak { a } }$ 计算公式，并计算室温下的 $K _ { \mathfrak { a } }$ 值。

表4.1 不同颜色的 NaHIn 溶液的吸光度 A/a.u

\textit{（原卷此处含表格/仪器试剂表，详见 PDF 或 MD 源文件。）}

\subsection*{4-2}
利用分光光度法测量配合物的稳定常数时，通过合理的实验设计，可以方便数据处理。

表 4.2 不同组成混合溶液的吸光度 A/a.u.

\textit{（原卷此处含表格/仪器试剂表，详见 PDF 或 MD 源文件。）}

用 $0 . 0 5 0 0 0 { \bmod { \mathrm { L } } } ^ { - 1 } { \mathsf { C u } } ( { \mathsf { N O } } _ { 3 } ) _ { 2 }$ (用M表示）和0.0500molL-1磺基水杨酸（用L表示）溶液配制系列混合溶液（将 pH调至4.5\~4.8后，定容至50.00mL），在440nm波长 $( \mathsf { C u } ( \mathsf { N O } _ { 3 } ) _ { 2 }$ 与磺基水杨酸的吸收可忽略）处测定各溶液的吸光度A，结果如表4.2所示。画出A-x图及必要的辅助线，计算并确定磺基水杨酸铜配合物的组成和表观稳定常数。

第6页共17页

第38届中国化学奥林匹克（决赛）理论第一场试题 2024年10月26日广州
